\documentclass[11pt,a4paper]{article}
\usepackage[T1]{fontenc}
\usepackage[utf8]{inputenc}
\usepackage[french]{babel}
\usepackage{lmodern}
\usepackage{geometry}
\geometry{margin=2.5cm}
\usepackage{amsmath,amssymb,booktabs,array}
\usepackage{graphicx}
\usepackage{hyperref}
\usepackage{tikz}
\usepackage{fancyhdr}
\usepackage{xcolor}
\hypersetup{colorlinks=true,linkcolor=blue!55!black,urlcolor=blue!55!black}
\pagestyle{fancy}
\fancyhf{}
\fancyhead[L]{Projet Data Science -- Commerce international}
\fancyfoot[C]{\thepage}
\newcommand{\achanger}[1]{\textcolor{blue!60!black}{\textbf{[À compléter : #1]}}}
\title{\textbf{Rapport de projet -- Data Science et commerce international}}
\author{\achanger{Binôme : Noms et prénoms}}
\date{\today}
\begin{document}
\begin{titlepage}
\centering
{\Large \achanger{Établissement} \par}
\vspace{1cm}
{\large Filière : \achanger{filière} \quad Module : Data Science\par}
\vspace{1.8cm}
{\Huge\bfseries \achanger{Titre précis du projet}\par}
\vspace{1.5cm}
{\Large Rapport scientifique de projet\par}
\vspace{1.5cm}
\begin{tabular}{rl}
Étudiant(e) 1 : & \achanger{Nom 1} \\
Étudiant(e) 2 : & \achanger{Nom 2} \\
Responsable : & A. Larhlimi \\
Année universitaire : & \achanger{2026--2027} \\
Date de remise : & \achanger{JJ/MM/AAAA} \\
\end{tabular}
\vfill
\textit{Document modèle : remplacer toutes les indications bleues. Ne jamais inventer des résultats.}
\end{titlepage}
\section*{Résumé}\addcontentsline{toc}{section}{Résumé}
\achanger{150 à 200 mots : contexte, jeu de données et URL, question de recherche, méthodes, résultats réellement obtenus, limites.}
\tableofcontents\newpage
\section{Introduction}\label{sec:intro}
\subsection{Contexte économique}\achanger{Décrire l'enjeu de commerce international avec sources.}
\subsection{Problématique et hypothèses}\achanger{Formuler la question et les hypothèses.}
\subsection{Objectifs}\achanger{Définir 2 ou 3 objectifs mesurables.}
\section{Présentation du jeu de données}\label{sec:data}
\subsection{Source, période et unité d'analyse}
\achanger{Nom officiel, organisme, URL, méthode d'extraction, période, dimensions effectivement observées.}
\subsection{Dictionnaire des variables}
\begin{table}[ht]\centering
\caption{Dictionnaire des variables -- à personnaliser}\label{tab:variables}
\begin{tabular}{lll}\toprule Variable & Unité / type & Rôle\\\midrule
\texttt{annee} & année & Index temporel\\
\texttt{valeur} & unité monétaire à préciser & Exemple de cible\\
\bottomrule\end{tabular}
\end{table}
\subsection{Qualité et préparation}\achanger{Manquants, doublons, changements de nomenclature, exclusions.}
\section{Analyse exploratoire}\label{sec:eda}
\achanger{Au moins cinq graphiques commentés dans le rapport ou son annexe.}
\begin{figure}[ht]\centering
% Remplacer le schéma d'exemple par un vrai fichier : \includegraphics[width=.8\linewidth]{figures/exportations.png}
\begin{tikzpicture}[>=stealth]
\draw[->] (0,0)--(9,0) node[right]{Temps};
\draw[->] (0,0)--(0,3) node[above]{Valeur};
\node[align=center,draw=blue!50!black,rounded corners,fill=blue!5] at (4.5,1.5)
{Insérer ici un graphique\\réel issu du notebook};
\end{tikzpicture}
\caption{Emplacement d'une figure expérimentale -- à remplacer}\label{fig:exemple}
\end{figure}
La figure~\ref{fig:exemple} devra être interprétée à partir des résultats calculés.
\section{Méthodologie}\label{sec:method}
\subsection{Prétraitement et protocoles}\achanger{Décrire pipeline, partage entraînement/validation/test et justification temporelle.}
\subsection{Modèles}\achanger{Présenter modèle de référence et deux modèles comparés.}
\subsection{Équations et notations}
Par exemple, l'erreur absolue moyenne (MAE) s'écrit :
\begin{equation}\label{eq:mae}
\mathrm{MAE}=\frac{1}{n}\sum_{i=1}^{n}\left|y_i-\widehat{y}_i\right|.
\end{equation}
Dans l'équation~\eqref{eq:mae}, $n$ est le nombre d'observations évaluées ; $y_i$ est la valeur réelle de la cible pour l'observation $i$ ; $\widehat{y}_i$ désigne la prédiction du modèle ; $|\cdot|$ est la valeur absolue. \textbf{Toute autre équation insérée dans le rapport doit définir chacun de ses paramètres.}
\section{Résultats expérimentaux}\label{sec:results}
\begin{table}[ht]\centering
\caption{Comparaison des performances -- ne pas renseigner de valeurs fictives}\label{tab:performances}
\begin{tabular}{lcc}\toprule Méthode & MAE & RMSE\\\midrule
Référence & \achanger{valeur} & \achanger{valeur}\\
Modèle 1 & \achanger{valeur} & \achanger{valeur}\\
Modèle 2 & \achanger{valeur} & \achanger{valeur}\\
\bottomrule\end{tabular}
\end{table}
Le tableau~\ref{tab:performances} sera adapté à la tâche (classification ou clustering si nécessaire).
\section{Discussion}\label{sec:discussion}
\achanger{Interpréter les résultats pour le commerce international ; biais, limites, restrictions causales.}
\section{Conclusion et perspectives}\label{sec:conclusion}
\achanger{Synthèse, réponse à la question de recherche, limites et extensions.}
\begin{thebibliography}{9}
\bibitem{comtrade} Nations Unies, \emph{UN Comtrade API -- Documentation}. \url{https://uncomtrade.org/docs/un-comtrade-api/}.
\bibitem{worldbank} Banque mondiale, \emph{Indicators API -- Documentation}. \url{https://datahelpdesk.worldbank.org/knowledgebase/articles/889392}.
\bibitem{wto} Organisation mondiale du commerce, \emph{Tariff and Trade Data}. \url{https://data.wto.org/en/dataset/wto_ttd}.
\end{thebibliography}
\appendix
\section{Annexe : reproductibilité}\achanger{Versions de Python et des bibliothèques, paramètres, répertoire des figures, sources et références complémentaires.}
\end{document}
